\section{Задание}

Смоделировать работу системы массового обслуживания на языке имитационного моделирования GPSS. Определить минимальный размер очереди, при котором не будет потерянных заявок. Генератор создаёт заявки, распределённые по равномерному закону на интервале [1, 5]. Обслуживающий аппарат обрабатывает заявки по нормальному закону с параметрами $\mu = 9$, $\sigma = 2$. Предусмотреть возможность возврата обработанной заявки в очередь с вероятностью 70\%. Смоделировать обработку 1000 заявок.

\section{Теоретическая часть}

\subsection{Равномерное распределение}

Случайная величина $X$ имеет равномерное распределение на отрезке $[a, b]$, если её плотность распределения $f(x)$ равна:

\begin{equation}
p(x) = \begin{cases}
\frac{1}{b-a}, & \text{если } a \leq x \leq b \\
0, & \text{иначе}
\end{cases}
\end{equation}

При этом функция распределения $F(x)$ равна:

\begin{equation}
F(x) = \begin{cases}
0, & x < a \\
\frac{x-a}{b-a}, & a \leq x \leq b \\
1, & x > b
\end{cases}
\end{equation}

Обозначение: $X \sim R[a, b]$.

Интервал времени между появлением $i$-й и $(i-1)$-й заявки по равномерному закону вычисляется следующим образом:

\begin{equation}
T_i = a + (b - a) \cdot R
\end{equation}

где $R$ — псевдослучайное число из интервала $[0, 1]$.

\subsection{Нормальное распределение}

Случайная величина $X$ имеет нормальное распределение с параметрами $\mu$ и $\sigma$, если её плотность распределения $f(x)$ равна:

\begin{equation}
f(x) = \frac{1}{\sigma \cdot \sqrt{2\pi}} e^{-\frac{(x-\mu)^2}{2\sigma^2}}, \quad x \in \mathbb{R}, \sigma > 0
\end{equation}

При этом функция распределения $F(x)$ выражается через функцию ошибок:

\begin{equation}
F(x) = \frac{1}{2} \cdot \left[ 1 + \text{erf}\left(\frac{x-\mu}{\sigma\sqrt{2}}\right) \right]
\end{equation}

Обозначение: $X \sim N(\mu, \sigma^2)$.

Интервал времени между появлением заявок по нормальному закону вычисляется с использованием центральной предельной теоремы:

\begin{equation}
T_i = \sigma \sqrt{\frac{12}{n}} \left( \sum_{j=1}^{n} R_j - \frac{n}{2} \right) + \mu
\end{equation}

При $n = 12$ формула упрощается:

\begin{equation}
T_i = \sigma \left( \sum_{j=1}^{12} R_j - 6 \right) + \mu
\end{equation}

где $R_j$ — псевдослучайные числа из интервала $[0, 1]$.

\subsection{GPSS}

Язык GPSS — общецелевая система моделирования.

Транзакты представляют собой описание динамических процессов в реальных системах. Они могут описывать как реальные физические объекты, так и нефизические. Транзакты можно генерировать и уничтожать в процессе моделирования.

Динамическими объектами являются транзакты, которые представляют собой единицы исследуемых потоков и производят ряд определённых действий, продвигаясь по фиксированной структуре.

Операционный объект. Блоки задают логику функционирования системы и определяют маршрут движения транзактов между объектами аппаратной категории.

К статическим объектам относятся очереди и таблицы, служащие для оценок влияющих характеристик.

Рассмотрим основные команды:

\begin{enumerate}
    \item \textbf{GENERATE} — команда, вводящая транзакты в модель
    \item \textbf{TERMINATE} — команда, удаляющая транзакт
    \item \textbf{QUEUE} — команда, помещающая транзакт в конец очереди
    \item \textbf{DEPART} — команда, удаляющая транзакт из очереди
    \item \textbf{SEIZE} — команда, занимающая канал обслуживания
    \item \textbf{RELEASE} — команда, освобождающая канал обслуживания
    \item \textbf{ADVANCE} — команда, задерживающая транзакт
    \item \textbf{TRANSFER} — команда, изменяющая движение транзакта в модели
    \item \textbf{START} — команда, управляющая процессом моделирования
\end{enumerate}

\section{Описание реализации}

\subsection{Структура программы}

В листинге 1 представлена реализация системы массового обслуживания на языке GPSS.

\includelistingpretty
    {lab6.gps}
    {gpss}
    {Реализация СМО на языке GPSS}

Программа состоит из следующих блоков:

\begin{itemize}
    \item \texttt{GENERATE (UNIFORM(1,1,5)),,,1000} — генератор создаёт 1000 заявок с интервалами, распределёнными равномерно на $[1, 5]$
    \item \texttt{QUEUE QSystemQueue} — заявка помещается в очередь
    \item \texttt{SEIZE Operator} — захват обслуживающего аппарата
    \item \texttt{DEPART QSystemQueue} — выход из очереди после захвата устройства
    \item \texttt{ADVANCE(NORMAL(1,9,2))} — обработка заявки по нормальному закону $N(9, 2)$
    \item \texttt{RELEASE Operator} — освобождение обслуживающего аппарата
    \item \texttt{TRANSFER 0.7,,Enqueue} — с вероятностью 70\% заявка возвращается в очередь, иначе завершается
    \item \texttt{TERMINATE 1} — удаление заявки и уменьшение счётчика завершения
    \item \texttt{START 1000} — запуск моделирования до обработки 1000 заявок
\end{itemize}

\section{Практическая часть}

В листинге 2 представлен результат работы программы.

\includelistingpretty
    {output.txt}
    {text}
    {Результат выполнения программы}

Основные результаты моделирования:

\begin{itemize}
    \item \textbf{Время моделирования:} 29493.663 единиц модельного времени
    \item \textbf{Сгенерировано заявок:} 1000
    \item \textbf{Всего входов в очередь:} 3269 (включая повторные обработки из-за возврата)
    \item \textbf{Максимальная длина очереди:} 896 заявок
    \item \textbf{Средняя длина очереди:} 451.830 заявок
    \item \textbf{Среднее время ожидания в очереди:} 4076.513 единиц времени
    \item \textbf{Утилизация обслуживающего аппарата:} 100\% (устройство ни разу не простаивало)
    \item \textbf{Среднее время обслуживания:} 9.021 (близко к заданному $\mu = 9$)
\end{itemize}

\textbf{Анализ результатов:}

Максимальная длина очереди составила \textbf{896 заявок}. Это означает, что для гарантированного обслуживания всех заявок без потерь необходимо обеспечить буфер очереди размером не менее 896 позиций.

Система работает в режиме перегрузки — обслуживающий аппарат загружен на 100\% времени. Это объясняется несколькими факторами:

\begin{enumerate}
    \item Высокая интенсивность генерации заявок (средний интервал 3 единицы времени)
    \item Относительно медленная обработка (среднее время 9 единиц)
    \item Эффект возврата заявок (70\% возвращаются в очередь), что значительно увеличивает нагрузку
\end{enumerate}

Из 1000 сгенерированных заявок система выполнила 3269 обработок (с учётом повторных). В среднем каждая заявка обрабатывалась $3269/1000 \approx 3.27$ раза.

\section{Вывод}

В ходе выполнения лабораторной работы была разработана модель системы массового обслуживания на языке GPSS с возможностью возврата обработанных заявок в очередь.

Программа корректно реализует заданную систему: генератор создаёт заявки по равномерному закону, обслуживающий аппарат обрабатывает их по нормальному закону, и 70\% заявок возвращаются для повторной обработки.

Определена минимальная длина очереди для обработки 1000 заявок без потерь — \textbf{896 позиций}. Система работает в режиме перегрузки с утилизацией обслуживающего аппарата 100\%, что подтверждает высокую нагрузку на систему. Среднее время ожидания в очереди (4076.513) значительно превышает среднее время обслуживания (9.021), что также свидетельствует о несбалансированности системы.

Для улучшения характеристик системы необходимо либо увеличить количество обслуживающих каналов, либо снизить вероятность возврата заявок, либо уменьшить интенсивность входного потока.
